% Dedication.
% STATUS: written 2026-11-06 on the owner's instruction: the dedication is
% the Arabic phrase "إلى المستقبل" (to the future), and the words are the
% owner's, not the machine's. Corrected the same day: the first render had the
% two words in the wrong order, see below.
%
% The book is typeset with pdfLaTeX, which cannot set Arabic. ArabTeX is
% installed but sets its transliteration font with \it, which memoir rejects,
% and it does not take Unicode input. Rather than change the engine of a
% 392-page book for one line, the phrase is set once by XeLaTeX with the Amiri
% font and stored as figures/dedication-arabic.png.
%
% The bidi package matters here. XeLaTeX with fontspec alone shapes the Arabic
% letters correctly but lays the run out left to right, so the words came out
% in reverse order: the file was unreadable to an Arabic reader while looking
% plausible to anyone checking only the glyph shapes. The page must be
% typeset inside an Arabic environment from polyglossia (or with bidi), which
% is what orders the words right to left.
%
% Regenerate it with:
%   article class, geometry paper 120mm x 30mm margin 0, polyglossia with
%   \setmainlanguage{english} and \setotherlanguage{arabic},
%   \newfontfamily\arabicfont[Script=Arabic]{Amiri}, and in a centred body:
%     {\fontsize{40}{48}\selectfont\arabicfont
%      \begin{Arabic}إلى المستقبل\end{Arabic}}
%   xelatex, then pdftoppm -png -r 300, then trim the white margin and save
%   as figures/dedication-arabic.png.
\thispagestyle{empty}
\vspace*{0.28\textheight}
\begin{center}
  \includegraphics[width=0.34\textwidth]{dedication-arabic.png}
\end{center}
\cleardoublepage
